
de présentation finie).

Ces propriétés étant locales sur @@M0@@ et @@M1@@, on peut les supposer affines; soient @@M2@@, @@M3@@. On a alors:
@@M4@@
d'après EGA IV 12.3.2 où on peut remplacer l'hypothèse de présentation finie sur @@M5@@ par des hypothèses noethériennes. Or @@M6@@ est isomorphe au corps résiduel de @@M7@@, et d'après EGA @@M8@@ 17.2.5 on a:
@@M9@@
Il résulte alors de EGA @@M10@@ 17.2.11 que cette dernière est finie si et seulement si @@M11@@ est de dimension cohomologique finie, c'est-à-dire régulier.

\textbf{Proposition 1.6.} \uline{Soient} @@M12@@ \uline{un schéma noethérien}, @@M13@@ \uline{un} @@M14@@\uline{-Module localement libre de type fini}, @@M15@@ \uline{le fibré vectoriel associé à} @@M16@@, \uline{et} @@M17@@ \uline{le morphisme structural de} @@M18@@. \uline{Alors} @@M19@@ \uline{est un isomorphisme}.

Comme @@M20@@ est lisse, et admet une section, à savoir la section nulle, il suffit pour prouver 1.6 de montrer que pour tout @@M21@@, @@M22@@ (1.4). On est ainsi ramené au cas où le schéma de base est un corps.

Dans ce cas, @@M23@@ est somme directe de sous-modules de rang @@M24@@, ce qui permet de décomposer @@M25@@ en une suite de fibrations de rang @@M26@@, et donc de se ramener au cas où @@M27@@ est de rang @@M28@@, soit @@M29@@. La proposition résulte alors du lemme suivant:

